% 第3章 方法

\section{整体架构}

本文的标注系统分三层，如图\ref{fig:architecture}所示。

% 图1: 三层系统架构
\begin{figure}[H]
\centering
\begin{tikzpicture}[node distance=0.8cm,
    box/.style={rectangle, draw=black, thick, minimum width=10cm, minimum height=1.2cm, align=center, font=\small}]
\node[box, fill=blue!8] (cl) {{\bf 数据合成层 -- Code2Logic}\\{游戏代码解析 $\to$ 模板填充 $\to$ 结构化 QA}};
\node[box, fill=green!8, below=of cl] (cg) {{\bf 场景理解层 -- Concept-Graphs}\\{SLAM $\to$ SAM $\to$ CF-SLAM $\to$ LLaVA $\to$ LLM $\to$ QA}};
\node[box, fill=orange!8, below=of cg] (is) {{\bf 真值验证层 -- Isaac Sim}\\{USD 场景生成 $\to$ 域随机化 $\to$ 几何真值导出}};
\node[box, fill=gray!10, below=of is] (out) {{\bf 输出: 结构化 QA 数据集 (2395 条, 5 种推理类型)}};
\draw[->, >=stealth, thick] (cl) -- (cg) node[midway, right] {\footnotesize 方法迁移};
\draw[->, >=stealth, thick] (cg) -- (is) node[midway, right] {\footnotesize 真值校准};
\draw[->, >=stealth, thick] (is) -- (out);
\end{tikzpicture}
\caption{系统三层架构。}
\label{fig:architecture}
\end{figure}


第一层是数据合成层。Code2Logic 从游戏源代码里提取逻辑结构，按模板生成 QA。这层的作用是证明方法论可行——不需要人工标注，从代码就能自动产出高质量问答数据。

第二层是场景理解层。把 Code2Logic 的模板化方法搬到真实 3D 场景。这里要处理的东西比游戏世界复杂得多——物体不是代码定义的坐标，而是从 RGB-D 数据里分割出来的；空间关系不是游戏规则写死的，要靠模型推断。

第三层是真值验证层。用 Isaac Sim 生成合成场景，每个物体的坐标、类别、尺寸都是精确已知的。以这些"标准答案"为基准，可以定量地算 Concept-Graphs 管线推断出来的标注到底有多准。

三层是递进关系：游戏方法在第 1 层被验证，搬到第 2 层解决真实问题，在第 3 层被校准。

\section{Concept-Graphs 标注管线}

\subsection{六步管线}

管线的数据流如图\ref{fig:pipeline}所示。

% 图2: 六步标注管线
\begin{figure}[H]
\centering
\begin{tikzpicture}[node distance=0.4cm,
    sbox/.style={rectangle, draw=black, thick, rounded corners=3pt, minimum width=1.8cm, minimum height=1.3cm, align=center, font=\footnotesize, fill=blue!5}]
\node[sbox] (s1) {RGB-D\\序列};
\node[sbox, right=of s1] (s2) {GradSLAM\\3D重建};
\node[sbox, right=of s2] (s3) {Grounded-\\SAM分割};
\node[sbox, right=of s3] (s4) {CF-SLAM\\3D建图};
\node[sbox, right=of s4] (s5) {LLaVA v1.5\\物体描述};
\node[sbox, right=of s5, fill=green!10] (s6) {llama.cpp\\GPU优化};
\node[sbox, right=of s6, fill=gray!15] (s7) {QA数据集};
\draw[->, >=stealth, thick] (s1) -- (s2);
\draw[->, >=stealth, thick] (s2) -- (s3);
\draw[->, >=stealth, thick] (s3) -- (s4);
\draw[->, >=stealth, thick] (s4) -- (s5);
\draw[->, >=stealth, thick] (s5) -- (s6);
\draw[->, >=stealth, thick] (s6) -- (s7);
\end{tikzpicture}
\caption{六步自动标注管线。}
\label{fig:pipeline}
\end{figure}


第一步是三维重建，用 GradSLAM\cite{gradslam} 从 RGB-D 序列和相机轨迹建点云地图，以 5 帧为步长处理约 400 个关键帧。第二步是二维分割，用 Grounded-SAM\cite{grounded_sam} 在每个关键帧上做检测和分割，这是最耗时的一步，单场景约 25 分钟。第三步是三维建图，用 CF-SLAM 把二维分割投影到三维空间，经重叠检测和 DBSCAN 聚类后输出每个物体的三维位置、包围盒和语义特征。

第四步是物体描述，用 LLaVA v1.5 7B\cite{llava} 对每个物体的裁剪区域生成自然语言描述。第五步是优化和推理，用 llama.cpp 在 GPU 上跑 qwen2.5:3b\cite{qwen2_5}，同时做两件事：优化 LLaVA 的原始描述（让它更像"物体标注"而非"场景描述"），推断物体之间的空间关系。第六步是按模板生成 QA 数据集。

\subsection{Prompt 设计}

LLaVA 输出的典型问题是两个：描述太笼统（"自然光洒满空间"），以及输出被截断（"es are all white..."）。针对这两个问题，本文设计了一套新的 prompt，核心改动有三条。

第一条是强制输出 class 字段——要求模型显式给出物体类别，例如 door、table、sofa。这直接解决了旧 prompt 输出"全是氛围、没有物体"的问题。第二条是设置禁用词汇表——light、atmosphere、depth、shadow 等十几个词被直接禁止。这是一个简单粗暴但极其有效的手段：模型生成这些词的概率被压到零，它就只能往"描述物体"的方向走。第三条是容错处理——输入太模糊时标记为 unknown，输入被截断时直接跳过。这避免了低质量数据流入下游。

Cheng 等人\cite{chain_of_thought} 的结论为这个设计提供了背书：对大模型来说，说清楚你要什么格式，比给一堆例子重要得多。

\subsection{GPU 加速}

原始管线用 Ollama 做本地推理，但 Ollama 内置的 CUDA 12.8 库（libggml-cuda.so）和服务器驱动（565.57.01，CUDA 12.7）不兼容。虽然 Ollama 声称把模型层"offload 到了 GPU"，但 nvidia-smi 上只显示 4 MiB 显存，实际计算全在 CPU 上——吞吐量只有约 3 tok/s，单场景处理要 15 分钟。

解决办法是换 llama-cpp-python。服务器上本来就装了 CUDA 12.6 完整工具链（nvcc 12.6、cmake 3.22、gcc 11.4），直接用这些系统库源码编译 v0.2.90。关键编译参数是 \texttt{-DGGML\_CUDA=on}，CUDA 架构设为 sm\_89。编译完成后，qwen2.5:3b GGUF 模型（Q4\_K\_M，1.9 GB）的 29 层全部加载到单张 RTX 4090 显存，开一个 OpenAI 兼容 API（端口 8080），管线里其它代码一行不用改。

编译中踩了三个坑。第一个是系统默认 nvcc 是 11.5，不支持 RTX 4090 的 sm\_89。通过 \texttt{-DCMAKE\_CUDA\_COMPILER} 指定 CUDA 12.6 的 nvcc 路径解决。第二个是缺少 OpenMP，用 \texttt{-DGGML\_OPENMP=off} 关掉。第三个是 CUDA 架构字符串格式——CMake 4.3 需要写成 \texttt{89} 而非 \texttt{sm\_89}。

优化后吞吐量 99 tok/s，显存占用 4.2 GB（7B 模型）或 1.9 GB（3B 模型）。

\section{Isaac Sim 验证系统}

\subsection{场景生成}

Isaac Sim 4.5 通过 pip 安装在服务器上，无头模式运行，Vulkan 渲染。场景生成完全程序化，不依赖图形界面。

设计了四类房型模板：起居室（6m×5m，10 个物体）、卧室（4.5m×4m，8 个物体）、厨房（5m×4m，8 个物体）、书房（4m×3.5m，7 个物体）。每类生成 5 个随机化变体，共 20 个场景。随机化覆盖三个维度：物体的位置在房间边界内随机扰动（碰撞检测保证间距 ≥0.5m），颜色从 3 种预设变体中随机选取，光照色温和强度连续随机采样。

\subsection{真值导出}

Isaac Sim 的 USD 场景图里，每个物体的坐标、类别和尺寸都是明确存储的，API 直接查询就能拿到标准答案，位置精度 0.01 米。空间关系用几何规则算：水平距离小于 1.5 米判定为 next to，1.5--3.0 米为 near，超出 3 米按方位角判断前后左右。因为计算过程是纯确定性的，准确率 100\%——不存在模型推断的不确定性。

场景图的导出格式和 Concept-Graphs 的标注模板一模一样，所以两边的数据可以直接对比。

\section{评估框架}

跨源评估从五个维度对比三类数据：（1）数据量——各源的 QA 总数和来源数；（2）QA 类型——Target Perception、Spatial Reasoning 等的分布；（3）关系类型——next to、above、behind 等的频次；（4）标注精度——类别准确率和关系准确率；（5）管线效率——数据获取方式和单场景耗时。
